% theme: multiple_and_half_angles
\MajorQuestionLesson{倍角・半角の公式}
\SetRowSpace{4mm}

\TypeQuestion{次の条件のもとで、$\sin2\theta,\ \cos2\theta,\ \tan2\theta$ の値を求めなさい。}{multiple_half_angles_double_value}
\QQRow{$\sin\theta=\dfrac{3}{5},\quad\dfrac{\pi}{2}<\theta<\pi$}{}{$\tan\theta=\dfrac{1}{3},\quad0<\theta<\dfrac{\pi}{2}$}{}
\QQRow{$\cos\theta=-\dfrac{5}{13},\quad\pi<\theta<\dfrac{3\pi}{2}$}{}{$\sin\theta=-\dfrac{2}{3},\quad\dfrac{3\pi}{2}<\theta<2\pi$}{}
\QQRow{$\tan\theta=\dfrac{1}{2},\quad\pi<\theta<\dfrac{3\pi}{2}$}{}{$\cos\theta=\dfrac{\sqrt{6}}{3},\quad0<\theta<\dfrac{\pi}{2}$}{}
\QQRow{$\tan\theta=-\dfrac{3}{2},\quad\dfrac{\pi}{2}<\theta<\pi$}{}{$\cos\theta=-\dfrac{1}{3},\quad\dfrac{\pi}{2}<\theta<\pi$}{}

\TypeQuestion{次の値を求めなさい。}{multiple_half_angles_half_exact}
\QQRowThree{$\sin\dfrac{\pi}{8}$}{}{$\tan\dfrac{3\pi}{8}$}{}{$\cos\dfrac{5\pi}{8}$}{}
\QQRowThree{$\cos\dfrac{\pi}{12}$}{}{$\sin\dfrac{7\pi}{8}$}{}{$\tan\dfrac{5\pi}{12}$}{}
\QQRowThree{$\tan\dfrac{7\pi}{8}$}{}{$\cos\dfrac{3\pi}{8}$}{}{$\sin\left(-\dfrac{\pi}{8}\right)$}{}

\TypeQuestion{次の条件のもとで、$\sin\dfrac{\theta}{2},\ \cos\dfrac{\theta}{2},\ \tan\dfrac{\theta}{2}$ の値を求めなさい。}{multiple_half_angles_half_value}
\QQRow{$\cos\theta=-\dfrac{7}{25},\quad\pi<\theta<\dfrac{3\pi}{2}$}{}{$\cos\theta=\dfrac{1}{3},\quad\dfrac{3\pi}{2}<\theta<2\pi$}{}
\QQRow{$\cos\theta=\dfrac{3}{5},\quad0<\theta<\dfrac{\pi}{2}$}{}{$\tan\theta=-\dfrac{3}{4},\quad\dfrac{\pi}{2}<\theta<\pi$}{}
\QQRow{$\cos\theta=-\dfrac{2}{3},\quad-\pi<\theta<-\dfrac{\pi}{2}$}{}{$\sin\theta=-\dfrac{12}{13},\quad\pi<\theta<\dfrac{3\pi}{2}$}{}

\LessonQuestionBreak

\TypeQuestion{次の式を簡単にしなさい。ただし、各式に含まれる分母は $0$ でないものとする。}{multiple_half_angles_simplify}
\QQRow{$2\sin\theta\cos\theta$}{}{$\dfrac{1-\cos2\theta}{2}$}{}
\QQRow{$\cos^2\theta-\sin^2\theta$}{}{$\dfrac{\sin2\theta}{1+\cos2\theta}$}{}
\QQRow{$\dfrac{\sin4\theta}{\sin2\theta}$}{}{$\dfrac{2\tan\theta}{1-\tan^2\theta}$}{}
\QQRow{$\cos^4\theta-\sin^4\theta$}{}{$\dfrac{1-\cos2\theta}{\sin2\theta}$}{}

\TypeQuestion{次の式を、三角関数の積や累乗を含まない形に直しなさい。}{multiple_half_angles_power_reduction}
\QQRow{$\sin^2\theta$}{}{$\cos^2 2\theta$}{}
\QQRow{$\sin^2\theta\cos^2\theta$}{}{$\sin^4\theta+\cos^4\theta$}{}

\TypeQuestion{次の方程式を、$0\leqq\theta<2\pi$ の範囲で解きなさい。}{multiple_half_angles_equation}
\QQRow{$\sin2\theta=\cos\theta$}{}{$\cos2\theta=\sin\theta$}{}
\QQRow{$\sin2\theta+\sin\theta=0$}{}{$\cos2\theta+\cos\theta=0$}{}
\QQRow{$\sin2\theta=2\sin\theta$}{}{$\cos2\theta=3\cos\theta-2$}{}
\QQ{$\tan2\theta+\tan\theta=0$}{}

\TypeQuestion{次の不等式を、$0\leqq\theta<2\pi$ の範囲で解きなさい。}{multiple_half_angles_inequality}
\QQRow{$\sin2\theta>\sin\theta$}{}{$\cos2\theta\leqq\cos\theta$}{}
\QQ{$\cos2\theta+\sin\theta\geqq0$}{}

\LessonQuestionBreak

\TypeQuestion{次の問いに答えなさい。}{multiple_half_angles_triple_identity}
\QQ{加法定理と倍角の公式を用いて、$\sin3\theta=3\sin\theta-4\sin^3\theta$ を証明しなさい。}{}
\QQ{加法定理と倍角の公式を用いて、$\cos3\theta=4\cos^3\theta-3\cos\theta$ を証明しなさい。}{}
\QQ{$4\sin^3 10^\circ-3\sin10^\circ$ の値を求めなさい。}{}
\QQ{$\cos^3\dfrac{2\pi}{9}-\dfrac{3}{4}\cos\dfrac{2\pi}{9}$ の値を求めなさい。}{}

\TypeQuestion{次の条件のもとで、$\sin3\theta,\ \cos3\theta$ の値を求めなさい。}{multiple_half_angles_triple_value}
\QQRow{$\sin\theta=\dfrac{3}{5},\quad\dfrac{\pi}{2}<\theta<\pi$}{}{$\cos\theta=-\dfrac{1}{3},\quad\pi<\theta<\dfrac{3\pi}{2}$}{}
\QQRow{$\tan\theta=\dfrac{1}{2},\quad0<\theta<\dfrac{\pi}{2}$}{}{$\sin\theta=-\dfrac{1}{2},\quad\dfrac{3\pi}{2}<\theta<2\pi$}{}
\QQRow{$\cos\theta=\dfrac{5}{13},\quad\dfrac{3\pi}{2}<\theta<2\pi$}{}{$\sin\theta=\dfrac{2}{3},\quad0<\theta<\dfrac{\pi}{2}$}{}

\TypeQuestion{次の方程式を、$0\leqq\theta<2\pi$ の範囲で解きなさい。}{multiple_half_angles_triple_equation}
\QQRow{$\sin3\theta=\sin\theta$}{}{$\cos3\theta=\cos\theta$}{}

\TypeQuestion{次の不等式を、$0\leqq\theta<2\pi$ の範囲で解きなさい。}{multiple_half_angles_triple_inequality}
\QQRow{$\sin3\theta-\sin\theta<0$}{}{$\cos3\theta+\cos\theta\geqq0$}{}
